% Note:
% At the moment 
% command format: 
% 1. \nomenclature[grouping-prefix]{symbol}{description}
% The grouping-prefix helps to sort the symbols: Symbols will be sorted by <prefox><symbol>. 

% defines the title of the list of symbols:
\renewcommand{\nomname}{List of Symbols}
% a short introduction to this section:
\renewcommand{\nompreamble}{Here lists several symbols that are used in this paper.}

% Grouping convention (Juat a first try, haven't figure out the best way to do it yet:
% c: small symbols
% g: Matrix and operators
% l: complex expressions.

\nomenclature[g]{$\hat{H}$}{Hamiltonian of the quantum mechanical system.}
\nomenclature[g]{$\hat{P_a}$}{Any Pauli operator or tensor of Pauli operators (Pauli string)}
\nomenclature[c]{$n$}{Number of basis functions used for in the Hamiltonian. It can also refer to, the number of fermionic modes, or the number of site on a lattice}
\nomenclature[c]{$m$}{The number of electrons in the system of interest}
\nomenclature[c]{$N$}{Refers to the number of qubits used to model the wavefunction of the system of interest}
\nomenclature[c]{$\epsilon$}{A given level of precision, i.e. the maximum error in estimation. A variant of this symbol, $\varepsilon$, is used to represent the level of noise in the quantum computer.}
\nomenclature[l]{$\mathcal{O}(f(n))$}{Order of function $f(n)$.}
\nomenclature[g]{$\unit$}{The identity matrix of an appropriate dimension.}
\nomenclature[g]{$I$}{The $2$ by $2$ identity matrix.}
\nomenclature[g]{$X/Y/Z$}{The three $2$ by $2$ Pauli $X$, Pauli-$Y$ and Pauli-$Z$ matrices.}
\nomenclature[g]{$\mathcal{P}$}{The number of Pauli strings in a Hamiltonian.}
\nomenclature[g]{$\hat{P}$}{Any Pauli string.}
\nomenclature[l]{$\phi_i(.)$, $\ket{\phi_i}$}{A spin-orbital function}
\nomenclature[l]{$\psi(.)$, $\ket{\psi}$}{An electronic wavefunction}

\printnomenclature
